\documentclass[10pt,a4paper]{amsart}
\usepackage[margin=19mm]{geometry}
\usepackage{amsmath,amssymb,mathtools}
\usepackage[scr=rsfs,cal=euler]{mathalfa}
\usepackage{xfrac}
\usepackage[unicode,hidelinks,bookmarks=false]{hyperref}
\newcommand{\ie}{i.e.\ }
\newcommand{\cf}{cf.\ }
\newcommand{\resp}{respectively\ }
\newcommand{\Subsection}{\subsection}
\providecommand{\nobreakdash}{\nobreak\hbox{-}}
\setlength{\emergencystretch}{2em}
\multlinegap0pt
\usepackage{bm}
\usepackage{bbm}
\usepackage{dsfont}
\usepackage{xspace}
\usepackage{cancel}
\usepackage{mathtools}
\usepackage{amsthm}
\makeatletter
\@ifundefined{theorem}{\newtheorem{theorem}{Theorem}}{}
\@ifundefined{definition}{\newtheorem{definition}[theorem]{Definition}}{}
\@ifundefined{lemma}{\newtheorem{lemma}[theorem]{Lemma}}{}
\makeatother
\title[The pairwise Stone space of an S4 De Morgan algebra]{The pairwise Stone space of an S4 De Morgan algebra}
\author{Joseph McDonald, Filip Jankovec}
\date{Source: arXiv:2608.03679v1 (CC BY 4.0)}
\begin{document}
\maketitle

\noindent\emph{Source and licence.} The pairwise Stone space of an S4 De Morgan algebra. Version arXiv:2608.03679v1 (CC BY 4.0), distributed under the Creative Commons Attribution 4.0 International licence, as recorded in the arXiv licence metadata for this exact version and shown on the abstract page https://arxiv.org/abs/2608.03679v1. Full rights notice: licenses/arxiv-2608.03679-CC-BY-4.0.txt.\par\smallskip

\noindent\textit{Editorial scope.} Counted unit: the Lemma whose source label is \texttt{homo to cont} in the article (source0.tex lines 548--550, source label \texttt{homo to cont}), delivered together with the article's own complete original proof of it (source0.tex lines 552--572, one \texttt{proof} environment of 21 lines). No mathematics is written, repaired or re-derived: statement and proof are the article's own text line for line. Required context retained uncounted: rep thm (lines 377--379); De Morgan homomorphism (lines 515--521); parwise frame homomorphism (lines 523--530). Every definition, display and label of those lines is retained unchanged. Omitted from this package and not redistributed: 1--376, 380--514, 522--522, 531--547, 551--551, 573--1007. Nothing inside a retained excerpt is omitted or altered: every retained body is delivered exactly as the article's lines are written, so no omission receipt of any kind accompanies this package. Citations: the retained excerpts carry no citation command at all, so no bibliography entry of the article is transcribed and no reference is redistributed. Reference adaptation: none was needed and none was made; every reference inside the delivered unit resolves inside it, so no unresolved reference or citation is produced. Printed slips kept literal: no printed word, symbol, hypothesis or number of the counted statement or of its proof was altered. Layout and class: the article's own class is \texttt{amsart} with packages including geometry, graphicx, amsfonts, amssymb, dirtytalk, color, amsmath, amsthm, multicol, mathtools, tikz, tikz-cd, tcolorbox, hyperref; the wrapper is the lane's pinned \texttt{amsart} editorial wrapper at 10pt on A4, which restates the theorem-like environments the retained text needs and adds the article's own macro definitions that the retained text needs (none), with extra wrapper packages (bm, bbm, dsfont, xspace, cancel, mathtools). The delivered numbering is the unit's own. Assets: no retained span contains a figure environment, a graphics-inclusion command, a PDF-page inclusion, a table or a TikZ picture; no image file is copied into this package, the package declares no asset dependency and no image file is redistributed with it. Licence: version arXiv:2608.03679v1 of this article is distributed under the Creative Commons Attribution 4.0 International licence, as recorded in the arXiv licence metadata for this exact version and shown on the abstract page https://arxiv.org/abs/2608.03679v1. A full rights notice is delivered at licenses/arxiv-2608.03679-CC-BY-4.0.txt. Characters: every retained excerpt is pure ASCII, so the delivered engine, XeLaTeX with its default Latin Modern text font, reports no missing character.
\begin{theorem}\label{rep thm}
    Every S4 De Morgan algebra $A$ is isomorphic to $\mathcal{A}_0(\mathcal{S}_0(A))$. 
\end{theorem}

\begin{definition} \label{De Morgan homomorphism}
    Let $A$ and $A'$ be S4 De Morgan algebras. A function $h\colon A\to A'$ is a \emph{homomorphism} provided: 
    \begin{enumerate}
        \item $h$ is a bounded lattice homomorphism; 
        \item $h(-a)=-h(a)$ and $h(\nabla a)=\nabla h(a)$. 
    \end{enumerate}
\end{definition}

\begin{definition}\label{parwise frame homomorphism}
    Let $X= \langle X, g, R, \tau_1, \tau_2 \rangle$ and $X' = \langle X', g', R', \tau'_1, \tau'_2 \rangle$ be PS4D-spaces. \\ A function $f\colon X\to X'$ is a \emph{pairwise frame morphism} provided:
    \begin{enumerate}
        \item $f$ is a bicontinuous function; 
        \item $f(g(x))=g'(f(x))$ for all $x\in X$; 
        \item $f^{-1}[R'[U]]=R[f^{-1}[U]]$ for all $U\in\tau'_1\cap\delta'_2$. 
    \end{enumerate}
\end{definition}

\begin{lemma}\label{homo to cont}
    Let $A$ and $B$ be S4 De Morgan algebras and let $h:A \to B$ be a homomorphism. Then $\mathcal{S}_1(h):\mathcal{S}_0(B) \to \mathcal{S}_0(A)$ defined by $\mathcal{S}_1(h)=h^{-1}$ is a bicontinuous frame morphism.
\end{lemma}

\begin{proof}
We start by showing that $\mathcal{S}_1(h)$ is bicontinuous. It is enough to verify this on the basis $\{\varphi_+(a):a\in A\}$ and $\{\varphi_-(a):a\in A\}$.
We have the following.
\begin{align*}
    \mathcal{S}_1(h)^{-1}[\varphi_+(a)]&=(h^{-1})^{-1}[\varphi_+(a)]=\{x \in \mathcal{S}_0(B) : h^{-1}[x] \in \varphi_+(a)\}\\&=\{x \in \mathcal{S}_0(B) : a \in h^{-1}[x]\}=\{x \in \mathcal{S}_0(B) : h(a) \in x\}=\varphi_+(h(a)).
\end{align*}
Similarly, we have
\begin{align*}
    \mathcal{S}_1(h)^{-1}[\varphi_-(a)]&=(h^{-1})^{-1}[\varphi_-(a)]=\{x \in \mathcal{S}_0(B) : h^{-1}[x] \in \varphi_-(a)\}\\&=\{x \in \mathcal{S}_0(B) : a \notin h^{-1}[x]\}=\{x \in \mathcal{S}_0(B) : h(a) \notin x\}=\varphi_-(h(a)).
\end{align*}
It remains to check the conditions (2) and (3) from Definition  \ref{parwise frame homomorphism}.
To show $\mathcal{S}_1(h)(g_B(x))=g_A(\mathcal{S}_1(h)(x))$ we have
$$\mathcal{S}_1(h)(g_B(x))=\mathcal{S}_1(h)[\{b \in B : -b \notin x\}]=h^{-1}[\{b \in B : -b \notin x\}]=\{a \in A : -h(a) \notin x\}.$$ Using Definition \ref{De Morgan homomorphism} (2) this is equal to 
$$\{a \in A : h(-a) \notin x\}=\{a \in A : -a \notin h^{-1}[x]\}=g_A(h^{-1}[x])=g_A(\mathcal{S}_{1}(h)(x)).$$

It remains to show $\mathcal{S}_1(h)^{-1}[R[U]]=R[\mathcal{S}_1(h)^{-1}[U]]$ for all $U \in \tau'_1 \cap \delta'_2$. Recall that $\mathcal{S}_1(h)^{-1}[\varphi_+(a)]=\varphi_+(h(a))$ and $\mathcal{S}_1(h)^{-1}[\varphi_-(a)]=\varphi_-(h(a))$. Additionally, we have $\nabla_R\varphi_+(a)=\varphi_+(\nabla a)$ and $\nabla_R\varphi_-(a)=\varphi_-(\nabla a)$ by Theorem \ref{rep thm}. Hence:  
 $$\mathcal{S}_1(h)^{-1}[R[\varphi_+(a)]]=\mathcal{S}_1(h)^{-1}[\nabla_R \varphi_+(a)]=\mathcal{S}_1(h)^{-1}[\varphi_+(\nabla (a))]=\varphi_+(h (\nabla a)).$$
By Definition \ref{De Morgan homomorphism} (2) this equals
$$\varphi_+ (\nabla (h(a)))=\nabla_R \varphi_+(h(a))=\nabla_R(h^{-1})^{-1}[\varphi_+(a)]=\nabla_R \mathcal{S}_1(h)^{-1}[\varphi_+(a)].$$
Analogously, it holds that $\mathcal{S}_1(h)^{-1}[\nabla_R \varphi_-(a)]=\nabla_R \mathcal{S}_1(h)^{-1}[\varphi_-(a)]$.  
\end{proof}

\end{document}
